\documentstyle[%


righttag%


,11pt%


,amssymb%


,russian%               remove in Eng.


,izvru%                 Set izven% in Eng.


% ,izvru-ks             for brief communications


]{amsart}





% 


\newcommand{\sign}{\operatorname{sign}}


\newcommand{\sat}{\operatorname{sat}}


\newcommand{\argmax}{\operatornamewithlimits{arg\,max}}


\newcommand{\tts}[1]{}





%\hoffset=-15mm%                  For Eng.


\setcounter{page}{78}


\god{2000}


\nomer{12~(463)} %                        Remove in Eng.


%\nomer{Vol.\,44, No.\,12}%               Set in Eng.


%\str{\pageref{begin}--\pageref{end}}%  Set in Eng.


\udk{517.977}





\author{\it В.А.\,Срочко, Е.И.\,Пудалова}


% NB:   ^^ remove in Eng.


\title{Методы нелокального улучшения допустимых управлений


в линейных задачах с запаздыванием}


\thanks{Работа выполнена при финансовой


поддержке Российского фонда фундаментальных исследований


(проект 99-01-400).}





\date{}


%\pagestyle{myheadings}%         Set in Eng.


\pagestyle{plain} %              Remove in Eng.


\begin{document}


%\thispagestyle{plain}%          Set in Eng.


\maketitle


\label{begin}





В настоящее время перспективным направлением теории вычислительных


методов оптимального управления является разработка итерационных


процедур со свойством нелокального улучшения допустимых управлений по


целевому функционалу. Это значит, что соответствующие методы не связаны


с операцией варьирования управлений, т.\,е. реализуются без


параметрического поиска, что является существенным фактором


повышения эффективности расчетов. Для обыкновенных


динамических систем такие методы построены и апробированы в классе


квадратичных задач оптимального управления \cite{71}, \cite{70}. Их


теоретической основой служат нелокальные аппроксимации целевого


функционала вместе с техникой его фазовой регуляризации.





В данной работе исследуются возможности такого подхода на уровне


управляемых систем с постоянным запаздыванием по состоянию.


Предварительный анализ показал, что построение и конструктивное


использование нелокальных формул приращения в квадратичных задачах с


запаздыванием встречает принципиальные трудности.


В определенной степени это является отражением тех проблем,


которые возникали в теории необходимых условий оптимальности второго


порядка для систем с запаздыванием в части дифференциального описания


сопряженной матричной функции \cite{49}. Поэтому нелокальный анализ


задач с запаздыванием проводится пока для линейных моделей: линейный


терминальный функционал, линейная по состоянию система с управлением в


коэффициентах. В этом классе построены точные формулы приращения


функционала, на основе которых конструируются две нелокальные процедуры


улучшения и объединяющий их метод приращений. Метод не вырождается,


вообще говоря, на особых управлениях, не использует параметрическоe


варьированиe и представляется наиболее эффективным средством численного


решения линейных задач с запаздыванием (экономичность, простота


реализации, потенциал улучшения). Для построения альтернативного метода


в билинейной задаче с выпуклым множеством допустимых управлений


проводится регуляризация функционала по управлению. В результате


получается метод проекций со


свойствами нелокального улучшения и сходимости по невязке принципа


максимума.





\section{Постановка задачи. Формула приращения функционала}





Рассмотрим задачу минимизации функционала


\begin{equation}


\Phi(u) =\langle c,\,x(t_1)\rangle \rightarrow\min,


\label{eq:z1}


\end{equation}


определенного на траекториях линейной по состоянию системы с постоянным


запаздыванием $h>0$


\begin{equation}


\begin{array}{c}


\Dot{x}(t) = A_1(u(t),t)x(t) + A_2(u(t),t)x(t-h) + b(u(t),t), \\[3mm]


x(t)=x^0(t),\ \;t\in [t_0-h,t_0],


\end{array}


\label{eq:z2}


\end{equation}


в классе допустимых управлений $V$, состоящем из


кусочно-непрерывных на $T=[t_0,t_1]$


вектор-функций $u(t)$ со значениями в заданном множестве $U\subset R^r$:


\begin{equation}


u(t)\in U,\ \;t\in T=[t_0,t_1].


\label{eq:z3}


\end{equation}





Пусть матричные функции $A_1(u,t)$ и $A_2(u,t)$ и вектор-функция


$b(u,t)$ в системе (\ref{eq:z2}) непрерывны по совокупности своих аргументов


на прямом произведении $U\times T$, множество $U$ компактно в $R^r$.


Предположим, что начальная функция $x^0(t)$ непрерывна на $[t_0-h,t_0)$.





Получим формулу приращения функционала в задаче (\ref{eq:z1})--(\ref{eq:z3}).


Введем обозначение для правой части системы (\ref{eq:z2})


$$


f(x,y,u,t)=A_1(u,t)x+A_2(u,t)y+b(u,t)


$$


и образуем функцию Понтрягина с сопряженной переменной $\psi\in R^n$


$$


H(\psi,x,y,u,t)=\langle \psi,\,f(x,y,u,t)\rangle \,.


$$





Пусть $u(t)$, $v(t)$ --- допустимые управления


в задаче (\ref{eq:z1})--(\ref{eq:z3}),


$x(t,u)$, $x(t,v)$ --- соответствующие им фазовые траектории. Для приращения


$\Delta x(t)=x(t,v)-x(t,u)$ справедливо уравнение


\begin{gather*}


\Delta\Dot{x}(t) =f(x(t,v),x(t-h,v),v(t),t)-f(x(t,u),x(t-h,u),u(t),t), \\


%


\Delta x(t)=0,\ \;t\in [t_0-h,t_0].


\end{gather*}


Отсюда


\begin{multline}


\Delta\Dot{x}(t) =\Delta_{v(t)} f(x(t,v),x(t-h,v),u(t),t)+\\


+(f(x(t,v),x(t-h,v),u(t),t)-f(x(t,u),x(t-h,u),u(t),t)),


\label{eq:z4}


\end{multline}


где символ $\Delta_{v(t)}$ означает частное приращение по управляющей


переменной на паре $u(t)$, $v(t)$.





Положим $x(t,v)=x(t,u)+\Delta x(t)$ и, учитывая линейность вектор-функции


$f$ по $x$ и $y$, представим уравнение (\ref{eq:z4}) в виде


\begin{multline*}


\Delta\Dot{x}(t)=f_x(x(t,u),x(t-h,u),u(t),t)\Delta x(t)+ \\


%


+f_y(x(t,u),x(t-h,u),u(t),t)\Delta x(t-h)


+\Delta_{v(t)}f(x(t,v),x(t-h,v),u(t),t).


\end{multline*}


Поскольку $f_x=A_1(u,t)$, $f_y=A_2(u,t)$, то получаем


систему в приращениях


\begin{equation}


\begin{array}{c}


\Delta\Dot{x}(t) = A_1(u(t),t)\Delta x(t) + A_2(u(t),t)\Delta x(t-h) +


\Delta_{v(t)} f(x(t,v),x(t-h,v),u(t),t), \\[3mm]


\Delta x(t)=0,\ \;t\in T_0=[t_0-h,t_0].


\end{array}


\label{eq:z5}


\end{equation}





Далее, рассмотрим приращение функционала (\ref{eq:z1}) на паре


управлений $u$ и $v$


$$


\Delta_v\Phi(u)=\langle c\,,\Delta x(t_1)\rangle .


$$





Пусть $\psi(t)$, $t\in T$, --- произвольная кусочно-дифференцируемая функция


с условием\linebreak $\psi(t_1)=-c$.


В силу уравнения (\ref{eq:z5}) найдем производную


\begin{multline*}


\frac{d}{dt}\langle \psi(t),\Delta x(t)\rangle


=\langle \Dot{\psi}(t),\Delta x(t)\rangle +


\langle A_1(u(t),t)^T\psi(t),\Delta x(t)\rangle + \\


%


+\langle A_2(u(t),t)^T\psi(t),\Delta x(t-h)\rangle +


\Delta_{v(t)}H(\psi(t),x(t,v),x(t-h,v),u(t),t)\,.


\end{multline*}


Тогда после интегрирования по $t\in T$


\begin{multline*}


\Delta_v\Phi(u)=-\int_{t_0}^{t_1}


\langle \Dot{\psi}(t),\Delta x(t)\rangle dt-


\int_{t_0}^{t_1}\langle A_1(u(t),t)^T\psi(t),\Delta x(t)\rangle dt- \\


%


-\int_{t_0}^{t_1}\langle A_2(u(t),t)^T\psi(t),\Delta x(t-h)\rangle dt-


\int_{t_0}^{t_1}\Delta_{v(t)}H(\psi(t),x(t,v),x(t-h,v),u(t),t)dt.


\end{multline*}


Преобразуем интеграл, содержащий $\Delta x(t-h)$ $(\tau=t-h)$,


$$


\int_{t_0}^{t_1}\langle A_2(u(t),t)^T\psi(t),\Delta x(t-h)\rangle dt=


\int\limits_{t_0-h}^{t_1-h}


\langle A_2(u(\tau+h),\tau+h)^T\psi(\tau+h),\Delta x(\tau)\rangle d\tau.


$$


Дополнительно предположим, что $\psi(t)=0$, $t\in(t_1,t_1+h]$. Тогда


$$


\int_{t_0}^{t_1}\langle A_2(u(t),t)^T\psi(t),\Delta x(t-h)\rangle dt=


\int\limits_{t_0}^{t_1}


\langle A_2(u(t+h),t+h)^T\psi(t+h),\Delta x(t)\rangle dt.


$$


Таким образом, приращение функционала имеет вид


\begin{multline*}


\Delta_v\Phi(u)=-\int_{t_0}^{t_1}


\langle \Dot{\psi}(t)+A_1(u(t),t)^T\psi(t),\Delta x(t)\rangle dt- \\


%


-\int_{t_0}^{t_1}


\langle A_2(u(t+h),t+h)^T\psi(t+h),\Delta x(t)\rangle dt-


\int_{t_0}^{t_1}\Delta_{v(t)}H(\psi(t),x(t,v),x(t-h,v),u(t),t)dt.


\end{multline*}





Определим вектор-функцию $\psi(t)=\psi(t,u)$, $t\in T$, с помощью


дифференциального уравнения


\begin{equation}


\Dot{\psi}(t)=-A_1(u(t),t)^T\psi(t)-A_2(u(t+h),t+h)^T\psi(t+h),\ \;t\in T,


\label{eq:z7}


\end{equation}


с начальным условием


\begin{equation}


\psi(t)=


\begin{cases}


-c,& t=t_1;\\


0,& t\in(t_1,t_1+h].


\end{cases}


\label{eq:z8}


\end{equation}


В результате формула приращения функционала (\ref{eq:z1}) принимает


окончательный вид


\begin{equation}


\Delta_v\Phi(u)=-\int_{t_0}^{t_1}


\Delta_{v(t)}H(\psi(t,u),x(t,v),x(t-h,v),u(t),t)dt.


\label{eq:z9}


\end{equation}


``Симметричный'' вариант формулы (\ref{eq:z9}) получается с помощью замены


$u\Leftrightarrow v$


\begin{equation}


\Delta_v\Phi(u)=-\int_{t_0}^{t_1}


\Delta_{v(t)}H(\psi(t,v),x(t,u),x(t-h,u),u(t),t)dt.


\label{eq:z10}


\end{equation}





Полученные формулы являются точными в том смысле,


что не содержат остаточных членов тех или иных разложений.





Подинтегральные выражения в (\ref{eq:z9}), (\ref{eq:z10}) суть частные


приращения по управлению функции Понтрягина для определенной


совокупности аргументов. Такое представление дает возможность


нелокального улучшения управления $u(t)$ через операцию на максимум


функции $H$, что является основой для последующих процедур и методов


улучшения.





Определим максимизирующее управление для функции $H$


\begin{equation}


u^*(\psi,x,y,t)=\argmax_{u\in U}


H(\psi,x,y,u,t),\ \;\psi,x,y\in R^n,\ \;t\in T.


\label{eq:z15}


\end{equation}





Предположим, что соотношение (\ref{eq:z15}) определяет вектор-функцию


$u^*(\psi,x,y,t)$,


которая является кусочно-непрерывной по своим аргументам, т.\,е. может иметь


конечное число поверхностей разрыва. Каждая такая поверхность


задается уравнением вида $g(\psi,x,y,t)=0$, и вектор-функция


$u^*(\psi,x,y,t)$ определяется неоднозначно только на поверхности разрыва.





Принцип максимума для управления $u(t)$, $t\in T$, в задаче


(\ref{eq:z1})--(\ref{eq:z3}) можно представить в виде


$$


u(t)=u^*(\psi(t,u),x(t,u),x(t-h,u),t),\ \;t\in T.


$$





На основе формул приращения (\ref{eq:z9}), (\ref{eq:z10}) нетрудно


получить достаточные условия оптимальности для управления $u(t)$ в


задаче (\ref{eq:z1})--(\ref{eq:z3}).





Согласно первой из формул (\ref{eq:z9}) неравенство


\begin{equation}


\Delta_{v(t)}H(\psi(t,u),x(t,v),x(t-h,v),u(t),t)


\leq 0\,,\;\;t\in T\,,\;v(t)\in U,


\label{eq:z11}


\end{equation}


является достаточным условием оптимальности управления $u(t)$, $t\in T$.





Введем множество достижимости системы (\ref{eq:z2}) в момент времени


$t\in T$ в заданном классе допустимых управлений


$D(t)=\{ x(t,u),\;u\in V\}$, $t\in T$.


Достаточным условием выполнения неравенства (\ref{eq:z11}) является


соотношение


$$


\Delta_v H(\psi(t,u),x,y,u(t),t)\leq 0,\ \;v\in U,\ \;t\in T,\ \;x\in


D(t),\ \;y\in D(t-h).


$$


Представим его в экстремальной форме


\begin{equation}


u(t)=\argmax_{v\in U}H(\psi(t,u),x,y,v,t),


\ \;t\in T,\ \;x\in D(t),\ \;y\in D(t-h).


\label{eq:z12}


\end{equation}


Это достаточное условие оптимальности управления $u(t)$ в задаче


(\ref{eq:z1})--(\ref{eq:z3}).


В частности, из (\ref{eq:z12}) при $x=x(t,u)$,


$y=x(t-h,u)$ получаем принцип максимума для управления $u(t)$.





Отметим, что в условии (\ref{eq:z12}) может фигурировать любая оценка по


включению для $D(t)$, т.\,е. множество $D:D(t)\subset D$, $t\in T$.





Согласно формуле приращения (\ref{eq:z10}) условие


оптимальности $\Delta_v\Phi(u)\geq 0$ обеспечивается неравенством


\begin{equation}


\Delta_{v(t)}H(\psi(t,v),x(t,u),x(t-h,u),u(t),t)


\leq 0,\ \;t\in T,\ \;v(t)\in U.


\label{eq:z13}


\end{equation}





Пусть $Q(t)=\{ \psi(t,u),\ u\in V\}$ --- множество достижимости 


сопряженной системы (\ref{eq:z7})--(\ref{eq:z8}) в момент $t\in T$.


Тогда неравенство (\ref{eq:z13}) обеспечивается условием


$$


\Delta_vH(\psi,x(t,u),x(t-h,u),u(t),t)\leq 0,\ \;v\in U,\ \;t\in


T,\ \;\psi\in Q(t).


$$


Таким образом, экстремальное соотношение


\begin{equation}


u(t)=\argmax_{v\in U}H(\psi,x(t,u),x(t-h,u),v,t),\ \;t\in


T,\ \;\psi\in Q(t),


\label{eq:z14}


\end{equation}


является достаточным условием оптимальности управления $u(t)$ в задаче


(\ref{eq:z1})--(\ref{eq:z3}). При\linebreak $\psi=\psi(t,u)$ получаем принцип


максимума для управления $u(t)$.





Выделим частный случай задачи (\ref{eq:z1})--(\ref{eq:z3}), когда матричные


функции $A_i(u,t)$, $i={1,2}$, не зависят от управления $u$:


$A_i(u,t)\equiv A_i(t)$. Тогда принцип максимума и условия (\ref{eq:z12}),


(\ref{eq:z14}) эквивалентны.


Действительно, в этом случае сопряженная система (\ref{eq:z7})--(\ref{eq:z8})


не зависит от управления и имеет единственное решение $\psi(t)$, т.\,е.


$Q(t)=\{\psi(t)\}$, $t\in T$. Это значит, что условие (\ref{eq:z14})


эквивалентно принципу максимума. Поскольку


$$


H=\langle \psi,A_1(t)x+A_2(t)y+b(u,t)\rangle\,,


$$


то услoвие (\ref{eq:z12}) также эквивалентно принципу максимума.





В рамках общей линейной задачи (\ref{eq:z1})--(\ref{eq:z3}) соотношения


(\ref{eq:z12}), (\ref{eq:z14}) могут служить основой для построения процедур


улучшения управления $u(t)$.





\begin{note}


Принцип максимума в рассматриваемой задаче не является


достаточным условием оптимальности.


\end{note}





\begin{note}


Традиционная формула приращения функционала в задаче


(\ref{eq:z1})--(\ref{eq:z3}) легко получается из (\ref{eq:z9})


и имеет вид \cite{22}


\begin{align*}


\Delta_v\Phi(u)&=-\int_{t_0}^{t_1}


\Delta_{v(t)}H(\psi(t,u),x(t,u),x(t-h,u),u(t),t)dt+\eta, \\


%


\eta&=-\int_{t_0}^{t_1}


(\langle \Delta_v H_x,\Delta x(t)\rangle


+\langle \Delta_v H_y,\Delta x(t-h)\rangle)dt.


\end{align*}


Эта формула лежит в основе доказательства принципа максимума и носит локальный


характер (не позволяет организовать нелокальное улучшение).


\end{note}





\section{Процедуры улучшения допустимых управлений}





Решим задачу улучшения для управления $u(t)$ (найти управление $v(t)$ с


условием\linebreak $\Phi(v)\leq\Phi(u)$) на основе формулы приращения


(\ref{eq:z9}).





{\it Первая процедура улучшения}:





по данному $u\in V$ найдем сопряженную траекторию $\psi(t,u)$, $t\in T$;





образуем максимизирующее управление


$$


v^*(x,y,t)=u^*(\psi(t,u),x,y,t),\ \;x,y\in R^n,\ \;t\in T;


$$





найдем решение $x(t)$, $t\in T$, фазовой системы


\begin{gather*}


\Dot{x}(t) = f(x(t),x(t-h),v^*(x(t),x(t-h),t),t), \\


x(t)=x^0(t),\ \;t\in [t_0-h,t_0],


%\label{eq:z16}


\end{gather*}


вместе с управлением $v(t)=v^*(x(t),x(t-h),t)$, $t\in T$.





Покажем, что новое управление $v(t)$ обеспечивает


свойство улучшения $\Phi(v)\leq\Phi(u)$.


Отметим, что траектория $x(t)$ определяется системой


\begin{gather*}


\Dot{x}(t) = A_1(v(t),t)x(t) + A_2(v(t),t)x(t-h) + b(v(t),t), \\


%


x(t)=x^0(t),\ \;t\in [t_0-h,t_0],


\end{gather*}


т.\,е. соответствует управлению $v(t)$. При этом $v(t)$ удовлетворяет условию


максимума


$$


v(t)=\argmax_{w\in U}H(\psi(t,u),x(t,v),x(t-h,v),w,t)\,.


$$


Следовательно,


$$


\Delta_{v(t)}H(\psi(t,u),x(t,v),x(t-h,v),u(t),t)\geq 0,\ \;t\in T.


$$


Отсюда в силу формулы (\ref{eq:z9}) заключаем, что $\Delta_v\Phi(u)\leq 0$,


т.\,е. первая процедура для любого $u\in V$ вырабатывает новое


управление $v\in V$ со свойством нестрогого улучшения. Равенство


$v(t)=u(t)$, $t\in T$, означает, что исходное управление $u(t)$ удовлетворяет


принципу максимума.





Трудоемкость реализации первой процедуры улучшения --- две задачи Коши


(для сопряженной и фазовой систем). При этом считается, что условие


максимума функции $H$ реализуется аналитически.





Отметим, что данная процедура позволяет, вообще говоря, улучшать


управления, удовлетворяющие принципу максимума.


\medskip





{\bf Пример 1.}


\begin{gather*}


\Phi(u)=x_2(1)\to\min, \\


\Dot{x}_1(t)=u(t)+x_1(t-h),\ \;\Dot{x}_2(t)=-u(t)x_1(t), \\


x(t)=0,\ \;t\in T_0=[-h,0],\ \;h=\frac{1}{2}, \\


|u(t)|\leq 1,\ \; t\in [0,1].


\end{gather*}





В данном случае $H(\psi,x,y,u,t)=\psi_1(u+y_1)-\psi_2ux_1$,


сопряженная система


\begin{gather*}


\Dot{\psi}_1(t)=\psi_2(t)u(t)-\psi_1(t+h),\ \;\Dot{\psi}_2(t)=0, \\


\psi_1(t)=0,\ \;t\in [1,\,1+h], \\


\psi_2(t)=


\begin{cases}


-1,& t=1;\\


0,& t\in(1,1+h].


\end{cases}


\end{gather*}


Максимизирующее управление $u^*(\psi,x,y,t)=\sign(\psi_1-\psi_2x_1)$. Понятно,


что $\psi_2(t)=-1$, $t\in[0,1]$.





Рассмотрим управление $u(t)=0$. Ему соответствуют траектории $x(t,u)=0$,


$\psi_1(t,u)=0$, $t\in T$. Данное управление яляется особым:


$H_u(\psi(t,u),x(t,u),x(t-h,u),t)=0$, т.\,е. удовлетворяет принципу


максимума с вырождением.





Применим первую процедуру улучшения.


Вспомогательное управление $v^*(x,t)=\sign x_1$.


Рассмотрим первое фазовое уравнение при $u=v^*$


$$


\Dot{x}_1(t)=\sign x_1(t)+x_1(t-h),\ \;x_1(t)=0,\ \;t\in T_0.


$$


Оно имеет особое решение $x_1(t)=0$ (с соответствующим управлением $u(t)=0$).


Кроме того, уравнение имеет еще два решения:


$$


x_1(t)=


\begin{cases}


\pm t, & t\in [0,h]; \\


\pm\frac{1}{8}(2t+1)^2, & t\in [h,1],


\end{cases}


$$


с порождающими управлениями $v(t)=\pm 1$. При этом выполняется свойство


строгого улучшения,


поскольку $\Phi(v)=-\int\limits_{0}^{1}|x_1(t)|dt$,


$\Phi(u)=0$.





Опишем ``симметричную'' процедуру улучшения, отправляясь от формулы


приращения (\ref{eq:z10}).





{\it Вторая процедура улучшения}:





по данному $u\in V$ найдем фазовую траекторию $x(t,u)$, $t\in T$;





сформируем экстремальное управление


$$


v^*(\psi,t)=u^*(\psi,x(t,u),x(t-h,u),t),\ \;\psi\in R^n,\ \;t\in T;


$$





найдем решение $\psi(t)$, $t\in T$, сопряженной системы


\begin{gather*}


\Dot{\psi}(t)=-A_1(v^*(\psi(t),t),t)^T\psi(t)-


A_2(v^*(\psi(t+h),t+h),t+h)^T\psi(t+h), \\


\psi(t)=


\begin{cases}


-c,& t=t_1;\\


0,& t\in(t_1,t_1+h],


\end{cases}


\end{gather*}


вместе с управлением $v(t)=v^*(\psi(t),t)$, $t\in T$.





В данном случае $\psi(t)=\psi(t,v)$, $t\in T$,


$$


\Delta_{v(t)}H(\psi(t,v),x(t,u),x(t-h,u),u(t),t)\geq 0,\ \;t\in T,


$$


и улучшение обеспечивается на основании формулы (\ref{eq:z10}).





Трудоемкость реализации второй процедуры улучшения --- две задачи Коши


(аналогично первой процедуре улучшения).


\medskip





{\bf Пример 2.}


\begin{gather*}


\Phi(u)=x_3(2)\to \min,\\


\Dot{x}_1(t)=u(t),\ \;\Dot{x}_2(t)=x_1(t-1),\\


\Dot{x}_3(t)=(x_2(t)+x_2(t-1)+x_3(t-1))u(t),\\


x_i(t)=0,\ \;t\in [-1,0],\ \;i=\overline{1,3},\\


|u(t)|\leq 1,\ \;t\in T=[0,2].


\end{gather*}


Здесь $H(\psi,x,y,u,t)=\psi_1 u+\psi_2 y_1+\psi_3(x_2+y_2+y_3)u$, сопряженная


система


\begin{gather*}


\Dot{\psi}_1(t)=-\psi_2(t+1),\\


\Dot{\psi}_2(t)=-\psi_3(t)u(t)-\psi_3(t+1)u(t+1),\\


\Dot{\psi}_3(t)=-\psi_3(t+1)u(t+1),\\


\psi_{1,2}(2)=0,\ \;\psi_3(2)=-1, \ \;


\psi_i(t)=0, \ \;t\in (2,3],\ \;i=\overline{1,3}.


\end{gather*}


Максимизирующее управление


$u^*(\psi,x,y,t)=\sign(\psi_1+\psi_3(x_2+y_2+y_3))$.





Рассмотрим управление $u(t)=0$ с траекториями $x_i(t,u){=}0$,


$i{=}\overline{1,3}$,


$\psi_{1,2}(t,u){=}0$, $\psi_3(t,u)={-}1$, $t\in T$.


Данное управление является особым, т.\,к.


$$


H_u(\psi(t,u),x(t,u),x(t-1,u),t)=0.


$$





Применим вторую процедуру улучшения. Вспомогательное управление


$v^*(\psi,t)=\sign\psi_1$.


На отрезке $[1,2]$\; $\psi_1(t)=0$ независимо от управления. Следовательно,


управление $v(t)$ для $t\in [1,2]$ можно выбрать


произвольным образом с точностью до ограничения $| v(t)|{\leq} 1$. Положим


$v(t){=}C$, $C\ne0$, $|C|\leq 1$. Соответствующие


сопряженные траектории $\psi_1(t,v)=0$, $\psi_2(t,v)=C(t-2)$,


$\psi_3(t,v)=-1$, $t\in [1,2]$. Тогда на отрезке $[0,1]$ появляется возможность


определить управление $v(t)=-\sign C$ и сопряженные траектории


\begin{gather*}


\psi_1(t)=-\frac{C}{2}(t-1)^2,\\


\psi_2(t)=\frac{| C|}{2}-Ct-2C+(t+1)\sign C,\\


\psi_3(t)=C(t-1)-1,\ \;t\in [0,1].


\end{gather*}


Таким образом, выходное управление второй процедуры имеет вид


$$


v(t,C)=


\begin{cases}


-\sign C, & t\in [0,1]; \\


C, & t\in [1,2].


\end{cases}


$$





Нетрудно подсчитать значение функционала $\Phi(v)=-|C|/6$.


Наилучший выбор $C=\pm 1$, что приводит к двум управлениям


$v(t)=v(t,\,\pm 1)$ со свойством строгого улучшения. Заметим, что полученные


управления удовлетворяют принципу максимума без вырождения.





\section{Метод приращений}





Рассмотрим линейную задачу (\ref{eq:z1})--(\ref{eq:z3}). Построим


итерационный метод последовательного улучшения допустимых управлений на


основе описанных процедур. Напомним, что максимизирующее управление $u^*$


определяется формулой (\ref{eq:z15}).





Пусть $u^0\in V$ --- начальное управление. Найдем соответствующую траекторию


$x^0(t)$ и подсчитаем значение функционала


$\Phi(u^0)=\langle c,x^0(t_1)\rangle$.





Пусть на $k$-й итерации имеется допустимая пара $(u^k(t),x^k(t))$ вместе


со значением функционала $\Phi(u^k)$, $k=0,1,\dotsc$


Найдем при $v^*(t)=u^*(\psi(t),x^k(t),x^k(t-h),t)$


решение $\psi^k(t)$ сопряженной системы 


\begin{gather*}


\Dot{\psi}(t)=-A_1(v^*(t),t)^T\psi(t)-A_2(v^*(t+h),t+h)^T\psi(t+h), \\


\psi(t)=


\begin{cases}


-c,& t=t_1;\\


0,& t\in(t_1,t_1+h],


\end{cases}


%\label{eq:z17}


\end{gather*}


и сформируем управление


$$


v^k(t)=u^*(\psi^k(t),x^k(t),x^k(t-h),t),\ \;t\in T.


$$


Найдем решение $x^{k+1}(t)$, $t\in T$, фазовой системы при


$w^*(t)=u^*(\psi^k(t),x(t),x(t-h),t)$


\begin{gather*}


\Dot{x}(t) = f(x(t),x(t-h),w^*(t),t), \\


x(t)=x^0(t),\ \;t\in [t_0-h,t_0],


%\label{eq:z18}


\end{gather*}


и образуем соответствующее управление


$$


u^{k+1}(t)=u^*(\psi^k(t),x^{k+1}(t),x^{k+1}(t-h),t),\ \;t\in T.


$$


Подсчитаем значение функционала $\Phi(u^{k+1})=\langle c,x^{k+1}(t_1)\rangle$.


Итерация закончена.





Прокомментируем метод. Управление $v^k(t)$ получено на основе $u^k(t)$


с помощью второй процедуры улучшения, т.\,е. $\Phi(v^k)\leq\Phi(u^k)$.


Отметим, что


$\psi^k(t)=\psi(t,v^k)$. Переход $v^k\Rightarrow u^{k+1}$ ---


реализация первой процедуры улучшения для 


управления $v^k$, т.\,е. $\Phi(u^{k+1})\leq\Phi(v^k)$.


Заметим, что значение $\Phi(v^k)$ фактически в методе не вычисляется.





В результате итерации получаем двойное улучшение по функционалу


$$


\Phi(u^{k+1})\leq\Phi(v^k)\leq\Phi(u^k),


$$


причем каждое улучшение дается ``ценой'' решения одной задачи Коши.





В этой связи отметим, что каждая процедура в отдельности обеспечивает одно


улучшение за счет двух задач Коши.





Поскольку функционал $\Phi(u)$ в задаче (\ref{eq:z1})--(\ref{eq:z3})


ограничен снизу (в силу компактности множества $U$ семейство фазовых


траекторий линейной системы (\ref{eq:z2}) ограничено), и последовательность


$\{\Phi(u^k)\}$ является невозрастающей, то сходимость метода приращений


описывается ``скромным'' соотношением $\Phi(u^k)-\Phi(u^{k+1})\to


0$, $k\to\infty$. При этом величина приращения


$\delta(u^k)=\Phi(u^k)-\Phi(u^{k+1})$ не является, вообще говоря, невязкой


принципа максимума для управления $u^k$, т.\,е. возможна ситуация, когда


$\delta(u^k)=0$, но управление $u^k$ не удовлетворяет принципу максимума.


Для устранения этих недостатков метода приращений будем


использовать процедуру регуляризации относительно приращения $\Delta u^k$.





\section {Билинейная задача. Метод проекций }





Рассмотрим билинейную задачу оптимального управления в следующей постановке:





минимизировать линейный (терминальный) функционал


\begin{equation}


\Phi(u)=\langle c,x(t_1)\rangle \to\min\,


\label{eq:z19}


\end{equation}


на траекториях билинейной системы


\begin{equation}


\begin{array}{c}


\Dot{x}(t)=A_1(u(t),t)x(t)+A_2(u(t),t)x(t-h)+b(u(t),t), \\[3mm]


x(t)=x^0(t),\ \;t\in T_0=[t_0-h,t_0],


\end{array}


\label{eq:z20}


\end{equation}


\begin{gather*}


A_i(u,t)=A_{i0}(t)+\sum_{j=1}^{r}u_j A_{ij}(t),\ \; i=1,2, \\


b(u,t)=B(t)u+с(t)


\end{gather*}


при ограничениях на управление


\begin{equation}


u(t)\in U,\ \;t\in T=[t_0,t_1].


\label{eq:z21}


\end{equation}





Предположим, что $U$ --- выпуклое замкнутое множество.





Пусть $(u^0(t),x(t,u^0))$, $t\in T$, --- допустимая пара в задаче


(\ref{eq:z19})--(\ref{eq:z21}). Введем вспомогательный функционал


\begin{equation}


\begin{split}


\Phi_\alpha(u,u^0)&=\Phi(u)+\alpha J(u,u^0),\ \;\alpha>0, \\


J(u,u^0)&=\frac{1}{2}\int_T\|u(t)-u^0(t)\|^2 dt.


\label{eq:z22}


\end{split}


\end{equation}


Функционал $\Phi_\alpha$ определяет процедуру регуляризации, связанную


с базовым управлением $u^0$, и использует только приращение $\Delta u^0(t)$.





Поставим задачу улучшения управления $u^0$ по функционалу


$\Phi_\alpha$: найти управление $v^\alpha\in V$ с условием


$\Phi_\alpha(v^\alpha,u^0)\leq\Phi_\alpha(u^0,u^0)$. При этом имеет место


оценка уменьшения исходного функционала $\Phi$


\begin{equation}


\Phi(v^\alpha)-\Phi(u^0)\leq-\alpha J(v^\alpha,u^0).


\label{eq:z23}


\end{equation}





Для решения задачи улучшения введем функцию Понтрягина для функционала


$\Phi_\alpha$


$$


H_\alpha(\psi,x,y,u,t)=\langle \psi,A_1(u,t)x+A_2(u,t)y+


b(u,t)\rangle -\frac{\alpha}{2}\|u-u^0\|^2


$$


вместе с сопряженной системой, которая не зависит от параметра $\alpha$


и управления $u^0$,


\begin{equation}


\begin{array}{c}


\Dot{\psi}(t)=-A_1(u(t),t)^T\psi(t)-A_2(u(t+h),t+h)^T\psi(t+h), \\[3mm]


\psi(t)=


\begin{cases}


-c,& t=t_1;\\


0,& t\in(t_1,t_1+h].


\end{cases}


\end{array}


\label{eq:z24}


\end{equation}


Отметим, что


$H_\alpha(\psi,x,y,u,t)=H(\psi,x,y,u,t)-\frac{\alpha}{2}\|u-u^0\|^2$, где


$H=H_0$ --- функция Понтрягина для задачи (\ref{eq:z19})--(\ref{eq:z21}).


В данном случае функция $H$ линейна по $u$.





Первая формула приращения функционала $\Phi_\alpha$ на управлениях


$u$, $u^0$ имеет вид


\begin{equation}


\Delta\Phi_\alpha(u,u^0)=-\int_T\Delta_{u(t)}


H_\alpha(\psi(t,u^0),x(t,u),x(t-h,u),u^0(t),t)dt.


\label{eq:z25}


\end{equation}





Представление (\ref{eq:z25}) дает возможность построить процедуру нелокального


улучшения для функционала $\Phi_\alpha$. Реализуем эту возможность через


операцию на максимум гамильтониана $H_\alpha$. Введем максимизирующее


управление


$$


u_\alpha^*(\psi,x,y,t)=\argmax_{u\in U}H_\alpha(\psi,x,y,u,t)=


\argmax_{u\in U}\bigg(\langle H_u(\psi,x,y,t),u\rangle


-\frac{\alpha}{2}\|u-u^0\|^2\bigg).


$$


Нетрудно видеть, что управление $u_\alpha^*$, $\alpha>0$, можно представить


в терминах операции проектирования


$$


u_\alpha^*(\psi,x,y,t)=


P_U\bigg(u^0(t)+\frac{1}{\alpha}H_u(\psi,x,y,t)\bigg).


$$


Здесь $P_U$ --- оператор проектирования на множество $U$ в евклидовой


метрике.





Отметим, что управление $u_\alpha^*(\psi,x,y,t)$ определяется однозначно:


проекция на выпуклое и замкнутое множество $U$ существует и единственна.


Кроме того, вектор-функция $u_\alpha^*(\psi,x,y,t)$ непрерывно зависит


от $\psi$, $x$, $y$ (оператор $P_U$ удовлетворяет условию Липшица,


производная $H_u(\psi,x,y,t)$ непрерывна по совокупности ($\psi,x,y$)).





Сформулируем принцип максимума, используя оператор проектирования:


{\it для оптимальности допустимого управления $u^0(t)$, $t\in T$, необходимо,


чтобы для любого $\alpha>0$ выполнялось соотношение}


$$


u^0(t)=P_U\bigg(u^0(t)+\frac{1}{\alpha}


H_u(\psi(t,u^0),x(t,u^0),x(t-h,u^0),t)\bigg),\ \; t\in T.


$$





Возьмем за основу формулу приращения (\ref{eq:z25}) и опишем {\it первую


процедуру улучшения} для управления $u^0(t)$:


\begin{itemize}


\item[] найдем решение $\psi(t,u^0)$, $t\in T$,


сопряженной системы (\ref{eq:z24}) при $u=u^0$ и сформируем управление


$v_\alpha^*(x,y,t)=u_\alpha^*(\psi(t,u^0),x,y,t)$;


\item[] найдем решение $x^\alpha(t)$ фазовой системы (\ref{eq:z20})


при $u(t)=v_\alpha^*(x(t),x(t-h),t)$ вместе с управлением


$v^\alpha(t)=v_\alpha^*(x_\alpha(t),x_\alpha(t-h),t)$.


\end{itemize}





Понятно, что $x^\alpha(t)=x(t,v^\alpha)$, $t\in T$, причем управление


$v^\alpha$ определяется соотношением


$$


v^\alpha(t)=\argmax_{u\in U}


H_\alpha(\psi(t,u^0),x^\alpha(t),x^\alpha(t-h),u,t)=


P_U\bigg(u^0(t)+\frac{1}{\alpha}


H_u(\psi(t,u^0),x^\alpha(t),x^\alpha(t-h),t)\bigg).


$$





Следовательно, на основании формулы (\ref{eq:z25}) при $u{=}v^\alpha$ имеет


место улучшение $\Delta\Phi_\alpha(v^\alpha,u^0){\leq}0$, т.\,е. управление


$v^\alpha$ для любого $\alpha>0$ обеспечивает невозрастание функционала


$\Phi(u)$ с оценкой (\ref{eq:z23}).


Из этой оценки следует, что в случае $v^\alpha\ne u^0$


свойство строгого улучшения гарантируется


для любого $\alpha>0$: $\Phi(v^\alpha)<\Phi(u^0)$. Равенство


$v^\alpha=u^0$, $t\in T$, означает, что управление $u^0(t)$ удовлетворяет


принципу максимума.





Таким образом, первая процедура (при любом $\alpha>0$) улучшает любое


допустимое управление $u^0(t)$, не удовлетворяющее принципу максимума в


задаче (\ref{eq:z19})--(\ref{eq:z21}).





Величина $\delta_\alpha(u^0)=\Phi(u^0)-\Phi(v^\alpha)$ может служить


невязкой принципа максимума для управления $u^0$:


$\delta_\alpha(u^0)=0\;\Leftrightarrow\;\Phi(u^0)=


\Phi(v^\alpha)\;\Leftrightarrow\;u^0=v^\alpha$.





Представим формулу (\ref{eq:z25}) в альтернативном виде


$$


\Delta\Phi_\alpha(u,u^0)=-\int_T


\Delta_{u(t)}H_\alpha(\psi(t,u),x(t,u^0),x(t-h,u^0),u^0(t),t)dt


%\label{eq:z28}


$$


и опишем {\it вторую процедуру улучшения}:


\begin{itemize}


\item[] образуем экстремальное управление


$v_\alpha^*(\psi,t)=u_\alpha^*(\psi,x(t,u^0),x(t-h,u^0),t)$;


\item[] найдем решение $\psi^\alpha(t)$


сопряженной системы (\ref{eq:z24}) при


$u(t)=v_\alpha^*(\psi(t),t)$ вместе с управлением


$v^\alpha(t)=v_\alpha^*(\psi^\alpha(t),t)$, $t\in T$.


\end{itemize}





В результате имеет место улучшение $\Delta\Phi_\alpha(v^\alpha,u^0)\leq 0$ с


оценкой уменьшения (\ref{eq:z23}) для функционала $\Phi(u)$.





Свойства данной процедуры аналогичны предыдущему случаю.





Таким образом, регуляризация (\ref{eq:z22}) порождает две независимые


процедуры строгого улучшения для любого допустимого управления, не


удовлетворяющего принципу максимума. При этом, однако, теряется свойство


улучшения для управлений, удовлетворяющих принципу максимума (в силу


единственности выходного управления $v^\alpha$).





Для иллюстрации последнего факта рассмотрим пример 1.


Применим первую процедуру к особому управлению $u^0(t)=0$. В данном


случае


$$


u_\alpha^*(\psi,x,t)=P_U\bigg(u^0(t)+\frac{1}{\alpha}(\psi_1+x_1)\bigg)


=\sat\bigg(u^0(t)+\frac{1}{\alpha}(\psi_1+x_1)\bigg),


$$


где


$$


\sat z=


\begin{cases}


\sign z,& | z| >1;\\


z,& | z|\leq 1.


\end{cases}


$$


Согласно схеме первой процедуры


$v_\alpha^*(x_1)=\sat\frac{x_1}{\alpha}$. Далее решается уравнение


$$


\Dot{x}_1(t)=\sat\frac{x_1(t)}{\alpha}+


x_1(t-h),\ \;x_1(t)=0,\ \;t\in[-h,0].


$$


Понятно, что для любого $\alpha>0$ единственным решением является


$x_1^\alpha(t)=0$ с порождающим управлением $v^\alpha(t)=0$. Исходное


управление переходит на выход процедуры без улучшения (других выходных


управлений нет).





Объединим описанные выше процедуры в метод последовательных приближений


(метод проекций).





Зафиксируем параметр регуляризации $\alpha>0$.


Пусть $k=0, 1,\dotsc$, ($u^k(t),x^k(t)$) --- допустимая пара в задаче


(\ref{eq:z19})--(\ref{eq:z21}). Выделим управление


$$


v^k(\psi,t)=u_\alpha^*(\psi,x^k(t),x^k(t-h),t)


$$


и найдем решение $\psi^k(t)$ сопряженной системы (\ref{eq:z24}) при


$u(t)=v^k(\psi(t),t)$.


Пусть $v^k(t)=v^k(\psi^k(t),t)$ --- промежуточное управление. Образуем


управление


$$


w^k(x,y,t)=u_\alpha^*(\psi^k(t),x,y,t)


$$


и найдем решение $x^{k+1}(t)$ фазовой системы (\ref{eq:z20})


при $u(t)=$ $w^k(x(t),x(t-h),t)$.





Пусть $u^{k+1}(t)=w^k(x^{k+1}(t),x^{k+1}(t-h),t)$ --- соответствующее


управление.





Пара ($u^{k+1}(t),x^{k+1}(t)$) завершает итерацию


метода.





Отметим, что $\Phi(u^{k+1})\leq\Phi(v^k)\leq\Phi(u^k)$, причем знак


равенства означает принцип максимума для соответствующего управления и может


служить критерием остановки.





Обсудим вопрос о сходимости метода. Введем неотрицательную величину


$\delta(u^k)=\Phi(u^k)-\Phi(u^{k+1})$. Это невязка принципа максимума


для управления $u^k$.





Сформулируем утверждение о сходимости: {\it если функционал $\Phi(u)$


ограничен снизу на множестве $V$, то метод проекций обладает свойством


сходимости по невязке принципа максимума


$(\delta(u^k)\to 0$, $k\to\infty)$.}





\begin{thebibliography}{9}





\bibitem{71}


Срочко В.А., Душутина С.Н., Пудалова Е.И. 


{\em Регуляризация принципа


максимума и методов улучшения в квадратичных задачах оптимального


управления} // Изв. вузов. Математика. -- 1998. -- \No\,12. -- С.\,82--92.


\bibitem{70}


Срочко В.А., Антоник В.Г. {\em Метод проекций в


линейно-квадратичных задачах


оптимального управления} // Журн. вычисл. матем. и матем.


физ. -- 1998. -- Т.\,38. -- \No\,4. -- С.\,564--572.


\bibitem{49}


Габасов Р., Кириллова Ф.М., Мансимов К.Б. {\em Необходимые условия


оптимальности второго порядка} (обзор). -- Минск.: Ин-т матем.


АН БССР, 1982. -- 48~с.


\bibitem{22}


Габасов Р., Кириллова Ф.М. {\em Качественная теория оптимальных


процессов}. -- М.: Наука, 1971. -- 508 с.





\end{thebibliography}





\vskip 2cm





\post{\it Иркутский государственный университет}{\it Поступила}


\post{}{28.04.2000}





\label{end}


\end{document}


